\documentclass[10pt]{article}

\usepackage[letterpaper,margin=0.82in]{geometry}
\usepackage{amsmath,amssymb,mathtools,bm}
\usepackage{microtype}
\usepackage{enumitem}
\usepackage[hidelinks]{hyperref}

\setlength{\parindent}{1.25em}
\setlength{\parskip}{0pt}
\setlist[itemize]{leftmargin=2.1em,itemsep=0.25em,topsep=0.35em}

\newcommand{\R}{\mathbb{R}}
\newcommand{\one}{\mathbf{1}}
\newcommand{\diag}{\operatorname{Diag}}
\newcommand{\prox}{\operatorname{prox}}
\newcommand{\ind}{\mathrm{I}}
\newcommand{\pos}[1]{\left[#1\right]_{+}}

\title{\vspace{-2.2em}Log-Barrier Primal--Dual Augmented Lagrangians\\
\large A short note on smooth multiplier updates and implicit complementarity}
\author{}
\date{}

\begin{document}
\maketitle
\vspace{-2.4em}

\begin{abstract}
This note replaces the nonnegative-orthant indicator in a standard slack-variable quadratic program by a logarithmic barrier and applies a primal--dual proximal augmented Lagrangian method. The regularized inequality saddle point admits a closed-form smooth multiplier and slack update. This update replaces the hard orthant projection of ordinary primal--dual augmented Lagrangian methods by exactly the paired retraction used by Arrizabalaga, Tracy, and Manchester~\cite{arrizabalaga2026single,arrizabalaga2026implicit}, up to the natural scaling of the dual proximal parameter. We also give a symmetric reduced Newton system and distinguish this barrier method from ordinary active-set PDAL methods such as ProxQP~\cite{bambade2022proxqp}.
\end{abstract}

\section{Convex QP preliminaries}

Following~\cite{arrizabalaga2026single,arrizabalaga2026implicit}, consider the convex quadratic program
\begin{equation}
\begin{aligned}
\underset{x\in\R^n}{\operatorname{minimize}}\quad
    &\frac12x^\top Qx+q^\top x
&\qquad\Longleftrightarrow\qquad
\underset{x\in\R^n,\,s\in\R^p}{\operatorname{minimize}}\quad
    &\frac12x^\top Qx+q^\top x\\
\text{subject to}\quad
    &Ax=b,\qquad Gx\leq h,
&
\text{subject to}\quad
    &Ax=b,\qquad Gx+s=h,\qquad s\geq0.
\end{aligned}
\end{equation}
Here $Q\in\mathbb{S}_+^n$, $q\in\R^n$, $A\in\R^{m\times n}$, $b\in\R^m$, $G\in\R^{p\times n}$, and $h\in\R^p$. The variables $y\in\R^m$ and $z\in\R^p$ denote the multipliers of $Ax=b$ and $Gx+s=h$, respectively. The KKT residuals and cone conditions are
\begin{equation}
\begin{aligned}
r_t&:=Qx+q+A^\top y+G^\top z=0,
&r_c&:=z\odot s=0,\\
r_e&:=Ax-b=0,
&r_i&:=Gx+s-h=0,
\qquad s,z\geq0.
\end{aligned}
\end{equation}
Equivalently, the slack formulation has objective
\begin{equation}
f(x)+\ind_{\R_+^p}(s),
\qquad
f(x):=\frac12x^\top Qx+q^\top x,
\end{equation}
subject to the two affine equalities in~(1). For a central-path parameter $\kappa>0$, replace the indicator by the extended-valued log barrier
\begin{equation}
\phi_\kappa(s):=
\begin{cases}
-\kappa\displaystyle\sum_{i=1}^p\log s_i,&s>0,\\
+\infty,&\text{otherwise}.
\end{cases}
\end{equation}
The stationarity condition of the resulting barrier problem with respect to $s$ is $z-\kappa s^{-1}=0$. Hence every barrier solution satisfies
\begin{equation}
z\odot s=\kappa\one,
\qquad s,z>0.
\end{equation}

\section{Barrier primal--dual augmented Lagrangian}

Let $(x^k,y^k,z^k)$ be the current primal--dual center. Following the proximal method of multipliers~\cite{rockafellar1976}, introduce a primal proximal regularization $\rho>0$ and dual proximal regularizations $\mu_e,\mu_i>0$. The barrier PDAL saddle function is
\begin{align}
\mathcal{M}_{\rho,\mu_e,\mu_i,\kappa}^k(x,s,y,z)
:={}&f(x)+\phi_\kappa(s)
  +y^\top(Ax-b)+z^\top(Gx+s-h)\notag\\
&+\frac{\rho}{2}\lVert x-x^k\rVert_2^2
  -\frac{\mu_e}{2}\lVert y-y^k\rVert_2^2
  -\frac{\mu_i}{2}\lVert z-z^k\rVert_2^2.
\end{align}
One proximal PDAL step computes, exactly or inexactly,
\begin{equation}
(x^{k+1},s^{k+1},y^{k+1},z^{k+1})
\approx
\underset{x,\,s>0}{\operatorname{argmin}}
\;\underset{y,z}{\operatorname{argmax}}
\mathcal{M}_{\rho,\mu_e,\mu_i,\kappa}^k(x,s,y,z).
\end{equation}
The corresponding optimality equations are
\begin{equation}
\begin{aligned}
Qx+q+A^\top y+G^\top z+\rho(x-x^k)&=0,\\
Ax-b+\mu_e(y^k-y)&=0,\\
Gx+s-h+\mu_i(z^k-z)&=0,\\
z-\kappa s^{-1}&=0.
\end{aligned}
\end{equation}
At a fixed point of the proximal iteration, the proximal terms vanish and~(8) reduces to the KKT system of the barrier QP. Sending $\kappa\downarrow0$ recovers the original complementarity conditions under the usual regularity assumptions.

\section{Closed-form smooth multiplier update}

For fixed $x$, the equality equation in~(8) gives
\begin{equation}
y=y^k+\frac{1}{\mu_e}(Ax-b).
\end{equation}
The inequality equations also admit a closed-form solution. Define
\begin{equation}
t:=z^k+\frac{1}{\mu_i}(Gx-h),
\qquad
\nu:=\frac{\kappa}{\mu_i}.
\end{equation}
Eliminating $s$ from the last two equations in~(8) gives, componentwise,
\begin{equation}
z_i^2-t_i z_i-\nu=0.
\end{equation}
Only the positive root is compatible with the barrier domain, so
\begin{equation}
\boxed{
z=b_\nu(t),
\qquad
s=\mu_i b_\nu(-t),
\qquad
b_\nu(a):=\frac{a+\sqrt{a^2+4\nu}}{2},
\qquad
\nu=\frac{\kappa}{\mu_i}.}
\end{equation}
The identity $b_\nu(a)b_\nu(-a)=\nu$ immediately gives
\begin{equation}
z\odot s=\mu_i\nu\one=\kappa\one.
\end{equation}
Equivalently, define the scaled multiplier $u:=\mu_i z$ and
\begin{equation}
\bar t:=\mu_i z^k+Gx-h,
\qquad
\bar\kappa:=\mu_i\kappa.
\end{equation}
Using the scaling identity $c\,b_\nu(a)=b_{c^2\nu}(ca)$ gives the paired form
\begin{equation}
\boxed{
u=b_{\bar\kappa}(\bar t),
\qquad
s=b_{\bar\kappa}(-\bar t),
\qquad
u\odot s=\bar\kappa\one.}
\end{equation}
Thus the log-barrier replacement has a particularly simple interpretation in ordinary PDAL. The hard multiplier projection
\begin{equation}
z=\pos{z^k+\frac{1}{\mu_i}(Gx-h)}
\end{equation}
is replaced by the smooth map
\begin{equation}
z=b_{\kappa/\mu_i}\!\left(z^k+\frac{1}{\mu_i}(Gx-h)\right).
\end{equation}
Since $b_\nu(a)\to[a]_+$ as $\nu\downarrow0$, ordinary PDAL is recovered in the zero-barrier limit. In proximal notation, $b_\nu=\prox_{\nu(-\log)}$; see, for example, Parikh and Boyd~\cite{parikh2014proximal}.

\section{Implicit complementarity and a symmetric Newton system}

The closed-form update is useful when the primal and dual blocks are alternated. A ProxQP-style implementation instead solves the coupled proximal subproblem. Introduce the implicit-complementarity coordinate $v\in\R^p$ and retract after every step using
\begin{equation}
z=b_\kappa(v),
\qquad
s=b_\kappa(-v),
\qquad
v=z-s.
\end{equation}
The proximal equations~(8) reduce to the smooth residual system
\begin{equation}
\begin{aligned}
R_t^k(x,y,v)
&:=Qx+q+A^\top y+G^\top b_\kappa(v)+\rho(x-x^k)=0,\\
R_e^k(x,y)
&:=Ax-b+\mu_e(y^k-y)=0,\\
R_i^k(x,v)
&:=Gx+b_\kappa(-v)-h+\mu_i\bigl(z^k-b_\kappa(v)\bigr)=0.
\end{aligned}
\end{equation}
Let
\begin{equation}
B:=\diag\!\left(b_\kappa'(v)\right),
\qquad
C:=\diag\!\left(b_\kappa'(-v)\right)=I-B,
\qquad
D:=C+\mu_iB,
\qquad
W:=CD^{-1}.
\end{equation}
The direct linearization of~(19) has the block derivatives $G^\top B$ and $-D$. Subtracting $G^\top W$ times the inequality row from the stationarity row gives the symmetric system
\begin{equation}
\boxed{
\begin{bmatrix}
Q+\rho I-G^\top WG&G^\top&A^\top\\
G&-D&0\\
A&0&-\mu_eI
\end{bmatrix}
\begin{bmatrix}
\Delta x\\ \Delta v\\ \Delta y
\end{bmatrix}
=-
\begin{bmatrix}
R_t^k-G^\top W R_i^k\\
R_i^k\\
R_e^k
\end{bmatrix}.}
\end{equation}
After computing a step length $\alpha$, update $(x,y,v)$ and retract:
\begin{equation}
x\leftarrow x+\alpha\Delta x,
\qquad
y\leftarrow y+\alpha\Delta y,
\qquad
v\leftarrow v+\alpha\Delta v,
\qquad
z,s\leftarrow b_\kappa(v),b_\kappa(-v).
\end{equation}
The slack and multiplier remain positive and satisfy perturbed complementarity after every completed step. Moreover,
\begin{equation}
0\prec B,C\prec I,
\qquad
\min(1,\mu_i)I\preceq D\preceq\max(1,\mu_i)I,
\qquad
0\preceq W\preceq I.
\end{equation}
Consequently, for fixed $\mu_i>0$, the complementarity-dependent blocks in~(21) remain bounded as $\kappa\downarrow0$. This removes the explicit slack-to-multiplier ratios that occur in standard primal--dual interior-point reductions. It does not, by itself, guarantee that the complete system is well conditioned; conditioning also depends on the problem data, degeneracy, and proximal parameters.

For numerical evaluation, the algebraically equivalent branch
\begin{equation}
b_\nu(a)=\frac{2\nu}{\sqrt{a^2+4\nu}-a}
\qquad (a<0)
\end{equation}
avoids cancellation in $a+\sqrt{a^2+4\nu}$.

\section{Relation to ordinary PDAL and scope}

The construction above is a smooth barrier analogue of the primal--dual proximal augmented Lagrangian method underlying ProxQP. Its main distinctions are:
\begin{itemize}
\item \textbf{Smoothness.} For fixed $\kappa>0$, the hard projection and discrete active-set changes are replaced by the smooth map $b_{\kappa/\mu_i}$.
\item \textbf{Relaxed complementarity.} The method solves a barrier QP satisfying $s^\top z=p\kappa$, rather than the original QP with exact complementarity.
\item \textbf{Linear algebra.} Ordinary PDAL can discard inactive inequality rows after active-set identification. Barrier PDAL generally retains every inequality, but its implicit-complementarity blocks remain bounded.
\item \textbf{Differentiation.} The fixed-$\kappa$ solution map is smooth and therefore convenient for automatic and implicit differentiation. The resulting derivative is the derivative of the relaxed problem and depends on $\kappa$.
\end{itemize}

For a fixed $\kappa>0$, the proximal iteration~(7) targets the saddle point of the barrier problem under the usual convexity and existence assumptions. To solve the original QP, one should place this iteration inside an outer schedule $\kappa_j\downarrow0$, warm-starting each barrier problem. Changing $\kappa$ at every PDAL iteration without controlling the proximal-subproblem accuracy is instead an inexact path-following heuristic and requires a separate convergence argument.

An alternative for differentiable optimization is to solve the original QP with ordinary PDAL, relax the converged solution to a chosen $\kappa_{\mathrm{relax}}>0$, and differentiate the relaxed system. This separates nominal solution accuracy from gradient smoothness, following the strategy of~\cite{arrizabalaga2026single}.

\begin{thebibliography}{9}

\bibitem{arrizabalaga2026single}
J.~Arrizabalaga, K.~Tracy, and Z.~Manchester.
\newblock A differentiable interior-point method in single precision.
\newblock \emph{arXiv preprint arXiv:2605.17913}, 2026.
\newblock \url{https://arxiv.org/abs/2605.17913}.

\bibitem{arrizabalaga2026implicit}
J.~Arrizabalaga and Z.~Manchester.
\newblock Implicit primal--dual interior-point methods for quadratic programming.
\newblock \emph{arXiv preprint arXiv:2604.00364}, 2026.

\bibitem{bambade2022proxqp}
A.~Bambade, S.~El-Kazdadi, A.~Taylor, and J.~Carpentier.
\newblock PROX-QP: Yet another quadratic programming solver for robotics and beyond.
\newblock In \emph{Robotics: Science and Systems}, 2022.

\bibitem{rockafellar1976}
R.~T. Rockafellar.
\newblock Augmented Lagrangians and applications of the proximal point algorithm in convex programming.
\newblock \emph{Mathematics of Operations Research}, 1(2):97--116, 1976.

\bibitem{parikh2014proximal}
N.~Parikh and S.~Boyd.
\newblock Proximal algorithms.
\newblock \emph{Foundations and Trends in Optimization}, 1(3):127--239, 2014.
\newblock \url{https://doi.org/10.1561/2400000003}.

\end{thebibliography}

\end{document}
